\section{Discussion}\label{sec:discussion}

\subsection{What the in-simulator gain measures}

The central finding is that, on real programming data, a policy's gain inside a KT simulator mainly measures how much learning the simulator's model \emph{believes} its problem choices produce, not how much a learner would learn. \Cref{prop:insim} makes this precise: when a KT model is both the simulator and the yardstick, the expected reported gain of any policy is its accumulated drift, and the empirical surprise in DKT-A was indeed negligible. Six learner models with equivalent forecasting accuracy implied very different learning: the explicit-learning models attributed small, homogeneous gains to practice, whereas the two DKTs attributed gains of several percentage points to single problems, in directions that a second random seed reproduced only weakly. This is an instance of underspecification \citep{damour2022underspecification}: the data constrain the next-step forecast, which all models match, but not the counterfactual consequences of practice sequences that learners never followed. It extends the observation that student models of similar accuracy can prescribe different amounts of practice \citep{lee2012individualizing,rollinson2015predictive} from mastery decisions to optimised sequencing policies, where the discrepancy is amplified rather than merely exposed. Reinforcement learning is an efficient search for exactly such sequences. The more capable optimiser (PPO) exploited the DKT reward most, reaching an in-simulator gain about twice that of DQN and a phantom gain of the same size in the no-learning world.

Two further observations sharpen the interpretation. First, for a given policy the tracker's reported gain varied little across the nine worlds---for PPO with the DKT reward between {{tg_min_ppo_kt}} and {{tg_max_ppo_kt}} pp, for random selection between {{tg_min_random}} and {{tg_max_random}} pp: its value is fixed mostly by which problems the policy selects, not by how the simulated learner responds. Second, along learners' actual continuations DKT's drift averaged out; its forecasts were volatile, as reported for DKT before \citep{yeung2018dktplus}, but not systematically biased there. The problem arises off the data distribution, where a policy compounds drift step after step. Predictive validation on logged sequences cannot detect this failure, which is why models that were indistinguishable in \cref{tab:forecast} could disagree about policies.

\subsection{Relation to model exploitation and reward over-optimisation}

The phantom gain is the educational counterpart of model exploitation in model-based RL and of reward over-optimisation: an agent optimised against a learned model finds the actions for which the model is most optimistic \citep{kurutach2018metrpo,janner2019trust,gao2023scaling,skalse2022reward}. KT--RL recommenders are particularly exposed because the same learned model provides dynamics, reward and evaluation, so an evaluation inside that model cannot reveal the exploitation. Remedies developed in model-based RL suggest testable designs for educational recommenders: penalising the disagreement of an ensemble of trackers, of which our re-seeded twin is a minimal instance \citep{kurutach2018metrpo,yu2020mopo}; restricting planning to short horizons or well-supported problems \citep{janner2019trust}; or constraining the learner model so that its drift is interpretable and monotone \citep{zhang2023cmkt}. We did not evaluate these remedies, and whether they would recover real learning gains is an open question.

\subsection{Implications for the evaluation of adaptive recommenders}

Our results do not show that the published systems we build on are wrong; we did not reproduce their tasks, rewards or datasets. They show that a design shared by many of them---training and scoring a policy with the same fitted KT model, or with models fitted to the same logs---can report large improvements for policies that teach nothing. We therefore suggest four inexpensive checks for simulator-based studies of exercise recommendation. (i) Report a \emph{negative control}: evaluate every policy in a world without learning in which the true gain is known to be zero, and report the tracker's claimed gain there. (ii) Report the \emph{drift decomposition} of \cref{eq:decomposition}, which requires only the tracker and exposes gains that the model expects from itself. (iii) Evaluate in \emph{independently fitted worlds} with different learning structures and learners, following the robust evaluation matrix \citep{doroudi2017rem}, including at least one world whose learning parameters are explicit. (iv) Re-fit the tracker with a different seed and report the agreement of implied learning. None of these checks needs student data beyond the logs already used. Because RL results vary across training seeds, such checks should be reported per seed or with seed-level uncertainty \citep{henderson2018matters,agarwal2021precipice}.

\subsection{What reinforcement learning can offer}

Reinforcement learning was not useless. When the training world had an explicit, interpretable learning structure (AFM), DQN and PPO learned policies that transferred to independently fitted AFM and PFA worlds and to a world with a different, assumed learning law, approaching a privileged one-step oracle while spending fewer submissions than random selection. This agrees with the observation of \citet{doroudi2019reward} that policies derived from theory-constrained models are more trustworthy than policies derived from unconstrained data-driven models. The evidence is nevertheless limited: the gains were modest---{{afm_adv_min}} to {{afm_adv_max}} pp of expected post-test success over random selection across ten problems in AFM-B and PFA-B---they remain simulator quantities, the latent worlds share the AFM world's structure, and the same agents did not beat random selection in DKT-B. Rules that target a moderate predicted success (0.7) were also robust, whereas an easy-to-hard curriculum and easiest-first selection, the ``too easy'' extreme of the zone-of-proximal-development argument in the course brief, were consistently the least effective in the latent worlds.

\subsection{Limitations}

All policy outcomes are simulated. The latent worlds are fitted to observational logs of one platform and one course family; their learning coefficients are associations adjusted for ability and problem difficulty, not causal effects, and the ZPD world is an assumption. Because the latent worlds share an additive logistic structure, the robustness of agents trained in the AFM world is conditional on that structure. The skill grouping is ours, and pointers, although prominent in the course brief, were represented by only one pool problem because few pointer problems reached the 5\% support threshold. Effort is measured by submissions, not time on task. We studied first-attempt outcomes, a ten-problem horizon and a single effort weight; longer horizons, other rewards and other pools may change magnitudes, and we did not run sensitivity analyses on these choices. Each agent was trained with three seeds, and intervals condition on these seeds and on the fitted models; they are not adjusted for multiplicity, although all pre-registered contrasts are reported, and the rank correlations are computed over policies that are not independent of each other. The data come from one university's online judge, so generalisation to other languages, course designs and learner populations is untested. Finally, we did not test any policy with students; the evidence concerns the validity of simulator-based evaluation and the relative robustness of policies across plausible worlds, not learning gains of real learners.
